\chapter{RESULTADOS}
\label{ch:resultados}

% TODO: redigir a partir das análises existentes, conferindo arquivo, versão, denominador e unidade de cada resultado.
% Separar diagnósticos históricos das seis rodadas finais; não exigir novas simulações ou notas humanas.
\section{RESULTADOS DOS DETECTORES INICIAIS}

\section{ANÁLISES AFETIVAS EXPLORATÓRIAS}

\section{AVALIAÇÃO QUANTITATIVA E QUALITATIVA DO STDI}
% Distinguir os 50 pares históricos, os 36 pares primários da auditoria e os 18 casos diagnósticos.
% Relatar comparação de métodos e sensibilidades documentadas, sem alegar calibração humana dos escores.

\section{COBERTURA TEMÁTICA E DIAGNÓSTICO DE ATALHOS}

\section{COBERTURA DAS RODADAS FINAIS}
% Seis rodadas, quatro árvores e 1.200 nós únicos por rodada; conclusão de execução não valida qualidade.
% Conferir estados de processamento e ausências sem contar projeções, cache ou notas como novas gerações.

\section{COMPARAÇÃO PAREADA DAS TRAJETÓRIAS E DOS GERADORES}
% Parear por entrada dentro de cada corpus, preservando prefixos compartilhados e método STDI/VAD.
% Comparar distância ao original, mudança incremental e soma cumulativa; não atribuir diferenças entre corpora só ao idioma.
% Usar somente estatísticas e intervalos sustentados pelos artefatos, com sua unidade de reamostragem.

\section{ANÁLISE QUALITATIVA DE CASOS}
